\documentclass{scrreprt}

\usepackage{pylab}

\begin{document}
	
	\myheader{Варіант 1}
	
	{\large
		\begin{enumerate}
			\item \textbf{Наступний день.} Дата деякого дня визначається трьома натуральними числами: $d$ (день), $m$ (місяць), $y$ (рік). Визначити дату наступного дня. Розглянути випадки високосного і невисокосного року.
			\item \textbf{Квадрат.} Точки $A$, $B$, $C$, $D$ на площині задаються своїми координатами. Визначити, чи є чотирикутник $ABCD$ квадратом.
		\end{enumerate}
	}
	
	
	\vspace{2cm}
	\myheader{Варіант 2}
	
	{\large
		\begin{enumerate}
			\item \textbf{Упертий Біллі.} Юний Білл Гейтс не любив уроки граматики в школі і якось вирішив втекти з уроку через вікно. Кватирка має форму прямокутника роз-мі-ром $m\times n$ сантиметрів. Визначити, чи зможе голова Білла Гейтса пролізти у кватирку, якщо голова є ідеальною кулею діаметром $d$ сантиметрів.
			\item \textbf{Прямокутник.} Точки $A$ і $C$ на площині задані своїми координатами і є про-ти\-леж\-ними\-ми вершинами прямокутника, сторони якого паралельні осям координат . Визначити периметр та площу частини прямокутника, розташованої в першій координатній чверті.
		\end{enumerate}
	}
	
	\newpage
	
	\myheader{Варіант 3}
	
	{\large
		\begin{enumerate}
			%\item {\bfseries Игра <<Угадай число 1.0>>.} Компьютер <<загадывает>> случайное целое число от 0 до 9. Игроку предоставляется три попытки угадать это число. В ходе <<диалога>> компьютер не даёт подсказок, а только проверяет пра\-виль\-ность догадки. В случае успеха игроку выдаётся сообщение о выигрыше, а после третьей неудачной попытки "--- о проигрыше.
			\item \textbf{Великий комбінатор.} Новому директору фірми «Роги і копита» знадобилася нова шафа. Допоможіть Остапу визначити, чи пройде шафа у дверний проріз, якщо шафа має форму прямокутного паралелепіпеда зі сторонами $x$, $y$, $z$, а проріз дверей "--- прямокутну форму $a\times b$.
			\item \textbf{Перетин відрізків.} На площині задані чотири точки $A$, $B$, $C$, $D$ такі, що кожен із відрізків $AB$ і $CD$ паралельний одній з осей координат. Перевірити, чи перетинаються відрізки $AB$ та $CD$.
		\end{enumerate}
	}
	\vspace{2cm}
	
	\myheader{Варіант 4}
	
	{\large
		\begin{enumerate}
			\item \textbf{Хмарочос.} У хмарочосі $N$ квартир і один під'їзд. На кожному поверсі три квартири. Ліфт може зупинятися лише на непарних поверхах. Пасажир входить до кабіни ліфта і набирає номер необхідної йому квартири $M$. Визначити, який поверх ліфт повинен доставити пасажира.
			\item \textbf{Об'єднання прямокутників.} Два прямокутники на площині задані таким чином:
			$$
			\begin{array}{c}
				P_1 = \left\{(x, y):\quad a_1\leqslant x\leqslant a_2,\quad b_1\leqslant y\leqslant b_2\right\},\\ 
				P_2 = \left\{(x, y):\quad c_1\leqslant x\leqslant c_2,\quad d_1\leqslant y\leqslant d_2\right\}.
			\end{array}
			$$
			Знайти площу {\itshape об'єднання} цих прямокутників.
		\end{enumerate}
	}
	
	\newpage
	
	\myheader{Варіант 5}
	
	{\large
		\begin{enumerate}
			\item \textbf{Щасливий квиток.} За заданим шестизначним номером квитка визначити, чи він є <<щасливим>>. Квиток вважається <<щасливим>>, якщо сума перших трьох його цифр дорівнює сумі трьох останніх, але сума не дорівнює 13.
			\item \textbf{Паралельність.} Точки $A$, $B$, $C$, $D$ на площині задані своїми координатами. Чи паралельні прямі $AB$ и $CD$?
		\end{enumerate}
	}
	
	\vspace{2cm}
	
	\myheader{Варіант 6}
	
	{\large
		\begin{enumerate}
			\item \textbf{Касовий апарат.} Скласти програму, яка запитує у користувача вартість по-купки і суму грошей, внесену покупцем. Якщо здача не потрібна, надрукувати на екрані <<Спасибі>>; якщо внесено більше грошей, ніж потрібно, надрукувати <<Візьміть здачу... грн>> (із зазначенням потрібної суми); якщо ж грошей не-дос-та-тньо, надрукувати <<Для покупки внесіть ще ... грн>> також із зазначенням необхідної суми.
			\item \textbf{Трикутник і кола.} На площині задані координати вершин трикутника $ABC$ і два кола з центрами $\left(a_1, b_1\right)$, $\left(a_2, b_2\right)$ і радіусами $ r_1$, $r_2$ відповідно. Чи лежить цей трикутник повністю в області перетину цих кіл?
		\end{enumerate}
	}
	
	\newpage
	
	\myheader{Варіант 7}
	
	{\large
		\begin{enumerate}
			\item \textbf{Турист.} Турист рухався $t_1$ годин зі швидкістю $v_1$, потім $t_2$ годин "--- зі швидкістю $v_2$ і $t_3$ годин "--- зі швидкістю $v_3$. Визначити, який час він подолав половину маршруту.
			\item \textbf{Точка та відрізок.} Точки $A$ та $B$ на площині задані своїми координатами. Чи вірно, що відрізок $AB$ не паралельний жодній з координатних осей?
		\end{enumerate}
	}
	
	\vspace{2cm}
	
	\myheader{Варіант 8}
	
	{\large
		\begin{enumerate}
			%\item \textbf{Счастливый билет.} По заданному шестизначному номеру билета определить, является ли он <<счастливым>>. Билет считается <<счастливым>>, если сумма первых трёх его цифр равна сумме трёх последних, при этом сумма не должна равняться 13.
			\item \textbf{Аліса та Боб.} Дані дати народження Аліси та Боба. Кожна з дат складається з трьох чисел: день, місяць та рік. Визначити, хто з них старший.
			\item \textbf{Прямокутний трикутник.} Дано три довільні числа $a$, $b$, $c$. Визначити, чи є трикутник із сторонами $a$, $b$ і $c$ прямокутним. Якщо такого трикутника немає, виводити повідомлення. Програма має працювати коректно для будь-якого набору із трьох чисел.
		\end{enumerate}
	}
	
	\newpage
	
	\myheader{Варіант 9}
	
	{\large
		\begin{enumerate}
			\item \textbf{Чемпіонат.} Відомі результати двох спортивних команд за підсумками чемпіонату: кiль\-кiсть очок, число забитих м'ячів, число пропущених м'ячів. Ви\-зна\-чи\-ти, яка з команд вище в~турнірній таблиці: При рівності очок спочатку враховується різниця забитих і пропущених м'ячів, а потім кількість забитих м'ячів Випадок, коли результати команд ідентичні, не розглядати.
			\item \textbf{Трикутник.} Дані координати вершин трикутника на площині. Визначити, чи пара-лель-на одна з його сторін осі абсцис.
		\end{enumerate}
	}
	
	\vspace{2cm}
	
	\myheader{Варіант 10}
	
	{\large
		\begin{enumerate}
			\item \textbf{Попередній день.} Дата деякого дня визначається трьома натуральними числами: $d$ (день), $m$ (місяць), $y$ (рік). Визначте дату попереднього дня. Рік не є високосним.
			\item \textbf{Перетин прямокутників.} Два прямокутники на площині задані таким чином:
			$$
			\begin{array}{c}
				P_1 = \left\{(x, y):\quad a_1\leqslant x\leqslant a_2,\quad b_1\leqslant y\leqslant b_2\right\},\\ 
				P_2 = \left\{(x, y):\quad c_1\leqslant x\leqslant c_2,\quad d_1\leqslant y\leqslant d_2\right\}.
			\end{array}
			$$
			Знайти площу {\itshape перетину} цих прямокутників. Якщо пе-ре\-тин порожнiй, видати нульові значення.
			
		\end{enumerate}
	}
	
	\newpage
	
	\myheader{Варіант 11}
	
	{\large
		\begin{enumerate}
			%\item \textbf{Телефонная компания.} Услуги телефонной сети оплачиваются по следующему правилу: за разговоры до $A$ минут в~месяц включительно платится фиксированная сумма $B$ грн, а разговоры сверх установленной нормы оплачиваются по тарифу $C$ грн за минуту. Пользователь вводит количество минут разговоров за прошедший месяц. Вычислить плату за пользование телефоном.
			\item \textbf{Дати.} Дана дата з трьох чисел (день, місяць та рік). Вивести <<yes>>, якщо така дата існує (наприклад, 12.02.1999 "--- <<yes>>, 22.13.2001 "--- <<no>>). Розглянути випадки високосного та невисокосного років.
			\item \textbf{Перпендикулярність.} Дві прямі на площині задані своїми рівняннями $a_1 x+b_1 y +c_1=0$ і $a_2 x+b_2 y +c_2=0$. Визначити, чи перпендикулярні вони.
		\end{enumerate}
	}
	
	
	\vspace{2cm}
	\myheader{Варіант 12}
	
	{\large
		\begin{enumerate}
			%\item {\bfseries Игра <<Угадай число 1.0>>.} Компьютер <<загадывает>> случайное целое число от 0 до 9. Игроку предоставляется три попытки угадать это число. В ходе <<диалога>> компьютер не даёт подсказок, а только проверяет пра\-виль\-ность догадки. В случае успеха игроку выдаётся сообщение о выигрыше, а после третьей неудачной попытки "--- о проигрыше.
			\item \textbf{Робот.} Робот може переміщатися в чотирьох напрямках («11» "--- північ, "12" "--- захід, "13" "-- південь, "14" "--- схід) і приймати три цифрові команди: 0 "--- продовжувати рух, 1 "--- поворот наліво, –1 "--- поворот направо. Дане число (11, 12, 13 або 14) "--- вихідний напрямок робота і ціле число $N $ (0, 1 або -1) "--- послана йому команда. Вивести напрямок робота після виконання отриманої команди (тобто північ, захід, південь або схід).
			\item \textbf{Прямокутник.} Точки $A$ і $C$ на площині задані своїми координатами і є про\-ти\-леж\-ни\-ми вершинами прямокутника, сторони якого паралельні осям координат. Визначити, скільки частин розбивається прямокутник осями координат. Знайти площу кожної частини.
		\end{enumerate}
	}
	
	\newpage
	
	\myheader{Варіант 13}
	
	{\large
		\begin{enumerate}
			\item \textbf{Земельна ділянка.} Чи можна на прямокутній ділянці землі розміром $a\times b$ метрів розмістити дві будівлі з прямокутними основами розмірами $p\times q$ і $r\times s$ метрів? Побудови можна розташовувати лише паралельно сторонам ділянки.
			\item \textbf{Прямокутний трикутник.} Точки $A$, $B$, $C$ на площині задані своїми координатами. Визначити, чи є трикутник $ABC$ прямокутним. Розглянути всі можливі випадки розташування точок (збіг будь-яких двох точок; всі точки лежать на одній прямій). У разі виродженого трикутника виводитиме відповідне повідомлення.
		\end{enumerate}
	}
	
	\vspace{2cm}
	
	\myheader{Варіант 14}
	
	{\large
		\begin{enumerate}
			\item \textbf{День тижня.} Дано натуральне число $k$ ($1\leqslant k\leqslant 365$). Визначити, яким днем тижня буде $k$-й день року, якщо 1 січня цього року - вівторок.
			\item \textbf{Ромб.} Точки $A$, $B$, $C$, $D$ на площині задаються своїми координатами. Визначити, чи є чотирикутник $ABCD$ ромбом.
		\end{enumerate}
	}
	
	\newpage
	
	\myheader{Варіант 15}
	
	{\large
		\begin{enumerate}
			\item \textbf{Вік.} Відомо натуральне число $n\leqslant 100$ "--- вік людини. Вивести фразу виду: <<Мені 21 рік>>, <<Мені 12 років>>, <<Мені 32 роки>>.
			\item \textbf{Два прямокутники.} Два прямокутники на площині задані таким чином:
			$$
			\begin{array}{c}
				P_1 = \left\{(x, y):\quad a_1\leqslant x\leqslant a_2,\quad b_1\leqslant y\leqslant b_2\right\},\\ 
				P_2 = \left\{(x, y):\quad c_1\leqslant x\leqslant c_2,\quad d_1\leqslant y\leqslant d_2\right\}.
			\end{array}
			$$
			Чи перетинаються вони? (Випадок однієї загальної точки також вважати перетином.)
		\end{enumerate}
	}
	
	\vspace{2cm}
	
	\myheader{Варіант 16}
	
	{\large
		\begin{enumerate}
			\item \textbf{Три ферзи.} На шахівниці стоять три ферзи (розташування кожної фігури за\-да\-ється парою на\-ту\-ральних\-них чисел від 1 до 8). Вивести пари ферзей, що загрожують один одному.
			
			{\itshape Ферзь <<б'є>> по горизонталі, вертикалі та діагоналі.}
			\item \textbf{Система рівнянь.} Дане число $a$. Знайти розв'язок системи рівнянь
			$$
			\begin{cases}
				y=x+1,\\
				y=|a(x-1)|.
			\end{cases}
			$$
		\end{enumerate}
	}
	
	\newpage
	
	\myheader{Варіант 17}
	
	{\large
		\begin{enumerate}
			\item \textbf{День тижня.} Дано натуральне число $k$ ($1\leqslant k\leqslant 365$). Визначити, яким днем тижня буде $k$-й день року, якщо 1 січня цього року - вівторок.
			\item \textbf{Ромб.} Точки $A$, $B$, $C$, $D$ на площині задаються своїми координатами. Визначити, чи є чотирикутник $ABCD$ ромбом.
		\end{enumerate}
	}
	
	
	\vspace{2cm}
	
	\myheader{Варіант 18}
	
	{\large
		\begin{enumerate}
			\item \textbf{Два ферзя.} На шахівниці стоять два ферзя (розташування кожної фігури задається парою на \-ту\-раль\-них чисел від 1 до 8). Чи загрожують вони один одному?
			
			{\itshape Ферзь <<б'є>> по горизонталі, вертикалі та діагоналі.}
			\item \textbf{Перетин відрізків.} На площині задані чотири точки $A$, $B$, $C$, $D$. Перевірити, чи перетинаються відрізки $AB$ та $CD$.
		\end{enumerate}
	}
	
	\newpage
	
	\myheader{Варіант 19}
	
	{\large
		\begin{enumerate}
			\item \textbf{Ферзь і король.} На шахівниці стоять білий король і чорний ферзь (розташування кожної фігури задається парою натуральних чисел від 1 до 8). Чи загрожує ферзь королю?
			
			{\itshape Ферзь <<б'є>> по горизонталі, вертикалі та діагоналі.}
			\item \textbf{Система рівнянь.} Дане число $a$. Знайти рішення системи рівнянь
			$$
			\begin{cases}
				y=\max\{1, x\},\\
				y=|a\cdot x|-1.
			\end{cases}
			$$
		\end{enumerate}
		
		
		\vspace{2cm}
		
		\myheader{Варіант 20}
		
		{\large
			\begin{enumerate}
				\item \textbf{День тижня.} Дано натуральне число $k$ ($1\leqslant k\leqslant 365$). Визначити, яким днем тижня буде $k$-й день року, якщо 1 січня цього року - пятниця.
				
				\item \textbf{Трикутник і точка.} Точки $A$, $B$, $C$, $M$ на площині задаються власними координатами. Визначити, чи точка $M$ всередині трикутника $ABC$.
			\end{enumerate}
		}
	\end{document}